% !TEX root = ../projeto-pesquisa.tex
% =====================================================================
%  METODOLOGIA
% =====================================================================
\secao{Metodologia}
A metodologia consiste em uma pesquisa de abordagem quantitativa e experimental,
estruturada como um estudo de campo controlado, com o objetivo de avaliar a
eficácia de técnicas computacionais na identificação de imagens geradas por
inteligência artificial, especialmente \textit{deepfakes}. A investigação
baseia-se na comparação entre a percepção visual humana e o uso de ferramentas de
ciência de dados, como leitura de metadados e análise de padrões de compressão.

A ação de extensão será realizada com estudantes do ensino médio de instituições
da região de Santo André, envolvendo aproximadamente 40 participantes. A proposta
adota uma abordagem dialógica, na qual os alunos atuarão de forma ativa na
construção do conhecimento, compartilhando percepções, hipóteses e dificuldades
ao longo das atividades, promovendo a troca de saberes entre a equipe do projeto
e a comunidade escolar.

O projeto será desenvolvido em etapas interdependentes. Inicialmente, será feita a
curadoria de 20 imagens digitais de alta resolução, sendo 10 autênticas e 10
geradas por inteligência artificial. A seleção considerará critérios que aumentem
a complexidade da avaliação, como texturas, iluminação, profundidade de campo e
ausência de artefatos evidentes. Também serão padronizados formato do arquivo,
resolução e nível de compressão, garantindo consistência metodológica no
experimento.

Paralelamente, será desenvolvido um protótipo de ferramenta digital, na forma de
um \textit{site} ou aplicação simples, com funcionalidades voltadas à verificação
de imagens. Entre elas, destacam-se a extração de metadados EXIF, a identificação
de \textit{softwares} de edição, a análise de padrões de compressão e a aplicação
da técnica de Análise de Nível de Erro (ELA), permitindo evidenciar
inconsistências que podem indicar manipulação digital. O desenvolvimento
priorizará a clareza dos resultados, de modo que os participantes compreendam os
indícios apresentados.

Na sequência, será realizada a formalização ética, com autorização da instituição
de ensino e aplicação do Termo de Consentimento de Participação, assegurando
participação voluntária e tratamento anônimo dos dados para fins acadêmicos.

A aplicação ocorrerá em ambiente escolar, iniciando-se com uma breve introdução
sobre inteligência artificial, \textit{deepfakes} e desinformação digital, em
linguagem acessível, sem interferir no desempenho dos participantes. Em seguida,
os estudantes serão divididos aleatoriamente em dois grupos iguais: o primeiro
grupo realizará a avaliação apenas com base na percepção visual, enquanto o outro
grupo contará com apoio da ferramenta desenvolvida. Cada participante
classificará as 20 imagens como reais ou artificiais dentro de um tempo definido,
com respostas registradas em formulário digital, garantindo padronização das
condições da atividade.

Após a atividade, será realizada uma discussão coletiva, na qual serão
apresentados exemplos e resultados preliminares. Os estudantes poderão refletir
sobre seus acertos e erros, compartilhar estratégias e discutir dificuldades,
estimulando o pensamento crítico sobre a confiabilidade de conteúdos digitais e
os limites da percepção humana.

Posteriormente, os dados serão submetidos à análise estatística, comparando as
taxas de acerto entre os grupos e identificando padrões de erro e níveis de
dificuldade das imagens. Isso permitirá avaliar a eficácia das técnicas
computacionais e suas limitações frente ao avanço dos modelos generativos. Por
fim, os resultados serão organizados em relatórios e produções acadêmicas, com
gráficos e análises interpretativas. Também será realizada uma devolutiva à
instituição participante, com orientações sobre identificação de conteúdos
manipulados, reforçando o caráter educativo da proposta.

